%======================================================================
\chapter{Introduction}
%======================================================================

Performance analysis in competitive sports has evolved from subjective observation to data-driven insights, yet the gap between raw footage and actionable intelligence remains formidable. This report presents Tobio, an AI-powered platform that bridges this divide by transforming volleyball game footage into comprehensive analytical data through computer vision and machine learning.

\section{Background}

Hours of game footage sit unwatched on team drives across the country. The problem is not a lack of data—it is a lack of time. Coaches and athletes know that reviewing gameplay is one of the most effective paths to improvement, but the reality is harsh: a single weekend tournament generates more video than anyone has time to analyze. For teams operating without the resources of professional organizations, this wealth of information becomes a burden rather than an asset.

Volleyball presents a particularly acute case of this challenge. The sport moves quickly, with rallies that demand split-second decisions and precise execution. Yet the very speed that makes volleyball exciting also makes manual analysis painfully slow. Finding specific plays, tracking individual performance, or identifying patterns requires scrubbing through hours of footage frame by frame.

\section{Motivation and Objectives}

The motivation for this project emerged from direct experience with this problem as a member of the Carleton University volleyball team. Game after game was recorded, uploaded, and ultimately forgotten—not because the footage lacked value, but because extracting that value required more time than anyone had to give.

Tobio was conceived as a solution to this persistent challenge. Rather than replacing human judgment, the system aims to automate the mechanical aspects of video analysis, freeing coaches and players to focus on strategy and improvement. The platform is designed around three core objectives:

\begin{enumerate}
    \item Automated spatial tracking of the court boundaries, ball trajectory, and player positions throughout each frame of gameplay.
    \item Event detection and classification to identify key moments such as serves, spikes, blocks, and scoring plays.
    \item Data synthesis into actionable statistics and visualizations, including performance metrics and spatial heat maps that reveal patterns invisible to real-time observation.
\end{enumerate}

The constraint of working from single-camera footage, rather than the multi-angle setups used by professional broadcasting, reflects the reality of most amateur and collegiate programs. This constraint shaped every technical decision in the project.

\section{Related Prior Work}

The technical foundation for this project was established through work on Whiplash, an experimental system with a radically different purpose but surprisingly similar challenges. Whiplash was an attempt to build a physical device capable of moving a computer mouse to aim and click on targets in video games with superhuman speed and accuracy—essentially, a robot that could play first-person shooters.

Leading the computer vision component for Whiplash meant solving a problem that, at its core, is identical to the one faced in sports analysis: detecting and tracking moving objects in real-time with minimal latency. The difference was one of stakes and scale. In Whiplash, a few milliseconds of delay or a few pixels of error meant the difference between hitting and missing a target. That project demanded optimization at every level—model selection, inference speed, tracking algorithms—to achieve performance that could compete with human reflexes.

The lessons from that work translated directly to Tobio. While the objects being tracked changed from game targets to volleyballs and players, the underlying technical challenges remained: how to extract meaningful spatial information from video streams fast enough to be useful, and how to do so reliably across varying lighting conditions, camera angles, and visual clutter. The difference is that Tobio does not need to operate in real-time; it processes recorded footage. This relaxation of the temporal constraint allowed for a shift in focus from raw speed to accuracy and completeness, though the optimization mindset from Whiplash continued to shape the design.
